\documentclass[11pt]{article}

\usepackage[a4paper,margin=2.5cm]{geometry}
\usepackage{fontspec}
\setmainfont{Latin Modern Roman}
\usepackage{booktabs}
\usepackage{longtable}
\usepackage{array}
\usepackage{xcolor}
\usepackage{listings}
\usepackage{hyperref}
\usepackage{parskip}
\usepackage{amsmath}

\hypersetup{
  colorlinks=true,
  linkcolor=blue!50!black,
  citecolor=blue!50!black,
  urlcolor=blue!50!black
}

\lstset{
  basicstyle=\ttfamily\footnotesize,
  breaklines=true,
  columns=fullflexible,
  keepspaces=true,
  frame=single,
  backgroundcolor=\color{gray!5},
  captionpos=b
}

\newcommand{\commit}[1]{\texttt{#1}}
\newcommand{\note}[1]{\textit{#1}}

\title{Optimizing the Concurrent Multiversion B-Tree (cMVBT):\\
       A Bottom-Up Tour of the System, its Tricks, and its OSIC Integration}
\author{Amir Tonta}
\date{August 2026}

\begin{document}
\maketitle

\begin{abstract}
This document collects, in one place, the optimizations applied to the \texttt{BatStore}
implementation over roughly two and a half years and $\sim$280 commits. It is organized
bottom-up: we start at the physical layout of a page on disk/in memory, work our way up
through concurrency control, structural maintenance (splits/merges), garbage collection,
and root indexing, and only then reach the transaction layer, where the centerpiece
optimization of this document lives --- the \emph{self-overwrite} trick for repeated
updates to the same key within one transaction. We then step back and look at two
cross-cutting, system-wide themes: \emph{sharding} (partitioning shared, contended
bookkeeping state) and the integration of \emph{OSIC} (Ordered Snapshot Instant Commit),
the durability/visibility model from LeanStore's snapshot-isolation paper. A running
theme throughout is a recurring engineering pattern: early designs patched races with
atomics and memory fences; later designs instead restructured the data so the race
became structurally impossible.
\end{abstract}

For a concise table-first inventory of every optimization, its evidence-backed gain,
reason, risk, and current or reverted status, see the standalone
\path{docs/optimization_catalog.tex}. This document remains the detailed narrative;
the catalog is the quick-reference index.

\tableofcontents

\section{Introduction and roadmap}

\texttt{cMVBT} is a concurrent, in-memory, multiversion B-tree that implements the
asymptotically optimal MVBT structure of Becker et al.\ \cite{becker1996asymptotically},
extended with a full multiversion concurrency control (MVCC) protocol described in
\cite{tonta2026multiversion}. On top of the tree itself, the system integrates
\emph{Ordered Snapshot Instant Commit} (OSIC), the durability/visibility scheme from
LeanStore's scalable snapshot-isolation work \cite{alhomssi2023scalable}, so that
transaction commits become visible instantly (in commit order) while the write-ahead log
(WAL) is hardened to disk lazily, in batches, out of the commit's critical path.

Reading a codebase's git history commit-by-commit gives you \emph{what changed}; it
rarely gives you \emph{why}, and it never gives you the shape of the whole system. This
document reconstructs both: for every subsystem below we explain the current design,
then the sequence of failed/superseded designs that preceded it and the concrete bug or
profiling number that motivated each change.

\paragraph{How to read this document.} It is ordered from the ground up, each section
building on the one before it:

\begin{enumerate}
  \item \textbf{Physical layer} (\S\ref{sec:physical}): how a page stores multiple
        versions of a key.
  \item \textbf{Concurrency control} (\S\ref{sec:cc}): how threads latch pages built
        from \S\ref{sec:physical} without blocking each other unnecessarily.
  \item \textbf{Structural maintenance} (\S\ref{sec:smo}): how pages from
        \S\ref{sec:physical}, latched per \S\ref{sec:cc}, split and merge as they fill
        up or empty out.
  \item \textbf{The self-overwrite trick} (\S\ref{sec:selfoverwrite}): the centerpiece
        optimization, sitting squarely on top of \S\ref{sec:physical} and
        \S\ref{sec:smo} --- how the transaction layer avoids ever handing the
        structural-maintenance code in \S\ref{sec:smo} a page it cannot split or merge.
  \item \textbf{Garbage collection} (\S\ref{sec:gc}): how pages retired by
        \S\ref{sec:smo} are reclaimed once no live snapshot needs them.
  \item \textbf{Root indexing} (\S\ref{sec:root}): how a reader finds the correct root
        for its snapshot in the first place.
  \item \textbf{Sharding} (\S\ref{sec:sharding}): a cross-cutting pattern applied to
        several of the subsystems above (GC, per-worker bookkeeping, and the database's
        own table layout) to remove contention on shared state.
  \item \textbf{OSIC integration} (\S\ref{sec:osic}): the transaction-commit and
        durability layer sitting on top of everything else, and how it differs from
        (and reuses) the ideas in \cite{alhomssi2023scalable}.
  \item \textbf{Timeline} (\S\ref{sec:timeline}): a chronological index of the commits
        discussed, for cross-reference.
\end{enumerate}

\section{Physical layer: pages and multiversion records}
\label{sec:physical}

Every leaf page (\path{src/mv_page_model/leaf_page.rs}) is a fixed-capacity,
\emph{append-only} array:
\begin{lstlisting}
record_data: [MaybeUninit<RecordPoint<Key, Payload>>; NUM_RECORDS]
\end{lstlisting}
plus a packed \path{len} counter tracking active vs.\ dead slots
(\path{leaf_page.rs:145-147}). \path{push_uncommitted} (\path{leaf_page.rs:163-180})
only ever writes at the next free slot and asserts \path{index < NUM_RECORDS}. Crucially,
\textbf{there is no in-place removal}: an obsolete, aborted, or deleted entry is only ever
\emph{marked} dead/invalid via \path{VersionInfo::delete}/\path{invalidate}
(\path{src/mv_record_model/version_info.rs:147-164}); it is never physically dropped until
the next structural-maintenance operation (\S\ref{sec:smo}) rebuilds the page from
scratch.

This means a single key can accumulate a \emph{physical chain} of several
\path{RecordPoint} tuples in one page --- inserts append a tuple, updates normally mark
the predecessor deleted and append its replacement, while a delete only marks the
existing tuple. Lookups scan backward (\path{rfind}) for the most recent live entry
(\path{leaf_page.rs:335-337}). Each tuple carries a \path{VersionInfo}
(\path{version_info.rs:37-40}): an \path{insertion_stamp: TxStamp} (a packed
\path{(worker\_id, ts\_start)}, constant for an entire transaction --- not per write, see
\S\ref{sec:selfoverwrite}) and an optional \path{deletion_stamp}. This packing is what
later makes ``was this the newest version, and did I, this very transaction, write it''
a single equality check rather than a heuristic.

Structural maintenance (\S\ref{sec:smo}) decides what to keep during a rebuild via
\path{record_survives_gc} (\path{src/mv_tree/smo.rs:736-744}): a record survives if it is
live, \emph{or} if its deleting transaction is still unresolved and may therefore abort,
which requires the predecessor to remain available for undelete
(\path{ctx.is_snapshot_live}). Older readers are not the reason for copying that record:
they route through the retired source block, and \path{live_min_snapshot} prevents that
block from being reused until those readers finish. An \emph{uncommitted} replacement's
fate is nevertheless undetermined until commit or abort, so it cannot be discarded.
The self-overwrite and reinsert-reuse paths in \S\ref{sec:selfoverwrite} bound how many
such unresolved tuples one transaction can create for one key.

The internal (index) page format underwent a parallel simplification, discussed together
with garbage collection in \S\ref{sec:gc}: an early design tracked obsolete child
pointers with an atomic ``obsolete mark'' bit mutated in place on an already-published,
concurrently-read slot; the current design (\commit{fe03d90}) instead derives liveness
from the fact that live sibling key intervals are always pairwise disjoint, so no slot is
ever mutated after publication (\path{InternalPage::live_mask()}).

\section{Concurrency control: from lock coupling to optimistic lock coupling}
\label{sec:cc}

\subsection{The protocols that were tried}

The lock-mode taxonomy predates the multiversion (\path{mv_}-prefixed) rewrite and still
lives in a sister crate, \texttt{CCBPlusTree}, whose \path{LockingStrategy} enum defines
six variants:

\begin{itemize}
  \item \textbf{MonoWriter} --- a single writer, unlimited readers, no real concurrency
        control at all.
  \item \textbf{LockCoupling (LC)} --- classic hand-over-hand exclusive latching down the
        tree.
  \item \textbf{ORWC} --- ``Optimistic Read-Write Coupling'': readers are free; writers
        use a level-bounded, optimistic-then-blocking coupling scheme with an
        ``optimistic upgrade'' path.
  \item \textbf{OLC} --- Optimistic Lock Coupling (Leis-style): readers are fully
        lock-free, relying only on version validation; writers use a versioned,
        seqlock-like latch.
  \item \textbf{LHL} (LightweightHybridLock) --- level-bounded: optimistic below a
        threshold depth, exclusive above it.
  \item \textbf{HL} (HybridLocking) --- bounded-retry hybrid read/write locking.
\end{itemize}

Only \textbf{OLC} survives in the current codebase: the write/read fast paths
(\path{src/mv_query/olc_query.rs}, \path{src/mv_query/query.rs},
\path{src/mv_query/rand_query.rs}) are hard-wired to OLC-style traversal
(\path{traversal_write_olc}, \path{traverse_read_key}, \path{try_retire},
\path{upgrade_write_lock}). MonoWriter, LockCoupling, LHL and HL were prototyped in
early 2024 (\commit{dcf3b2b}) but never carried into the multiversion rewrite.

\subsection{The rise and fall of ORWC}

ORWC is worth its own subsection because it is the one protocol that was implemented,
iterated on for roughly a year, and then deliberately removed --- a useful cautionary
tale about optimistic upgrade paths.

\begin{itemize}
  \item \commit{d565b2d} (2024-01-22) --- first sign of trouble: \emph{``ORWC seem to
        have an issue.''}
  \item \commit{ac006fe} (2024-01-26) --- the root is made upgradeable optimistically
        ``by all protocols --- deadlock-free, including for ORWC'', adding a
        post-upgrade re-validation (\path{root_guard.deref().unwrap().len() == curr_len})
        because a naive before/after version compare cannot detect a writer already
        mid-flight at both sample points.
  \item \commit{48c4ded} (2024-05-06) --- \emph{``Updated OLC for actual required
        validations. Fixed ORWC.''} Redundant \path{fence(Acquire)}/\path{fence(SeqCst)}
        calls are stripped from the traversal, since OLC's own atomic load/CAS already
        carries the needed ordering.
  \item \commit{b1e7092} (2024-05-07) --- \emph{``Further fixed ORWC issue with upgrades
        and free readers sometimes encountering segfaults.''} Guard variants are reworked
        from holding raw pointers to holding a cloned \path{SmartCell} handle, because a
        raw pointer captured at borrow time could dangle across a concurrent reuse.
  \item \commit{f368155} (2024-05-07, same day) --- \emph{``Make standard ORWC mut
        borrow blocking and only the upgrade non-blocking.''} ORWC's mutable-borrow path
        had to give up non-blocking semantics entirely to stay correct; only the
        optimistic \emph{upgrade} attempt remained non-blocking.
  \item \commit{3c6da7b} (2024-12-06) --- \emph{``Deprecated ORWC.''} The README changes
        from ``ORWC: Unlimited all.'' to ``ORWC: Deprecated; Optimistic Upgrade not
        working.''
  \item \commit{5c26abb} (2024-12-08) --- \emph{``Purged ORWC $\to$ duplicate with OLC.
        Cleaned code accordingly.''} $\sim$460 lines removed.
\end{itemize}

\textbf{Root cause, in one sentence:} ORWC's optimistic upgrade path never reached a
provably correct state --- four separate fix attempts over a year (segfaults on free
readers, upgrade races, blocking/non-blocking inconsistency) --- and once OLC was fully
validated it gave the identical guarantee through one simpler seqlock mechanism, making
ORWC a pure maintenance liability rather than a genuinely distinct protocol worth
keeping.

\subsection{The SmartCell latch}

\path{src/mv_sync/smart_cell.rs} implements the versioned latch (``SmartCell'') that
every node uses:

\begin{itemize}
  \item \textbf{Layout}: \path{OptCell<E> { cell: SafeCell<E>, cell_version: AtomicVersion }}
        --- a single 64-bit versioned word per node, with three logical bits packed into
        its top nibble: a \path{WRITE_FLAG_VERSION} (seqlock ``locked'' bit), pin-related
        flags, and a \emph{permanent} \path{RETIRED_FLAG_VERSION} tombstone bit.
  \item \textbf{Optimistic read} (\path{borrow_read}): one \path{Acquire} load of
        \path{cell_version}, masking out the write/retired bits --- no lock taken.
  \item \textbf{Upgrade to exclusive} (\path{upgrade_write_lock}): a CAS from the
        captured read version to \path{WRITE_FLAG_VERSION | read_version}, failing fast
        if the retired bit is already set (a retired cell can never be write-locked
        again).
  \item \textbf{Unlock}: stores \path{(write_version+1) \^{} WRITE_FLAG_VERSION | retired_bit}
        with \path{Release}. Reading the current retired bit with a plain \path{Relaxed}
        load here is safe (nothing else can touch \path{cell_version} while the writer
        guard is alive) and, crucially, preserves a same-critical-section
        \path{mark_retired()} call instead of clobbering it (see \S\ref{sec:gc} for the
        bug this fixes).
  \item \textbf{\path{try_retire}}: a single CAS straight from a \path{Reader} snapshot
        to \path{RETIRED_FLAG_VERSION | read_latch}, skipping the intermediate exclusive
        state --- sound only for guards that, once excluded, are unconditionally consumed
        with no ``back out'' path.
\end{itemize}

\subsection{Chasing fences out of the hot path}

A second, quieter optimization story runs in parallel to the protocol story above: once
OLC's own atomic operations were trusted to carry the necessary ordering, standalone
memory fences bolted on earlier (out of caution) were incrementally identified as
redundant and removed. Table~\ref{tab:fences} lists the sequence.

\begin{longtable}{@{}p{2cm}p{2cm}p{5.2cm}p{5.3cm}@{}}
\toprule
\textbf{Commit} & \textbf{Date} & \textbf{Change} & \textbf{Rationale} \\
\midrule
\endhead
\commit{48c4ded} & 2024-05-06 & Removed \path{fence(Acquire)}/\path{fence(SeqCst)} in the
  traversal loop & The surrounding OLC load/CAS already provides the ordering. \\
\commit{3c6da7b} & 2024-12-06 & Removed paired \path{fence(Acquire)} duplicating
  \path{len.load(Acquire)} in page length accessors & An atomic load's own ordering
  already suffices. \\
\commit{a7bb29c} & 2025-04-01 & \emph{``Avoid unnecessary CAS failures by moving to
  acquire/release semantics.''} CAS failure ordering \path{Relaxed}$\to$\path{Acquire};
  unlock stores \path{Relaxed}$\to$\path{Release} & Establishes the happens-before edges
  OLC readers actually rely on, cutting spurious CAS retries from
  under-synchronized reads racing writers. \\
\commit{a518eed} & 2025-04-25 & \emph{``Updated atomics, removed fences.''} Deleted
  standalone fences immediately adjacent to an atomic store/load on the same word & Pure
  dead weight once the adjacent atomic op's ordering subsumed it. \\
\commit{2a2be93} & 2026-06-22 & Page soft-commit length bumps \path{Release}$\to$
  \path{Relaxed} & The subsequent write-unlock's own \path{Release} on the latch already
  publishes everything written before it in program order. \\
\commit{40b00fd} & 2026-06-30 & \path{Node::m_type} \path{AtomicUsize}$\to$\path{usize};
  removed \path{fence(Acquire)} in \path{Drop for Node} & The type tag is only ever
  mutated while the caller holds the exclusive write lock; the latch's own release
  already publishes it. \\
\commit{45c7c21}/\commit{fd92f38} & 2026-07-21 & ``Remove conservative child latch'':
  \path{try_retire} lets read-only, unconditionally-consumed guards skip a full
  write-lock upgrade during split/merge & Removes one CAS/lock-acquire per split/merge
  for guards with no backout path. \\
\bottomrule
\caption{Chronology of fence removal and ordering tightening in the OLC latch path.}
\label{tab:fences}
\end{longtable}

Two further, structural memory-layout optimizations sit alongside this fence-removal
story:

\begin{itemize}
  \item \commit{7f2bd09} (2026-05-21) --- \emph{``Pin Root* for root\_live reorg\ldots
        Move node tag to single cacheline.''} \path{Node} gains
        \path{\#[repr(C, align(64))]} plus explicit padding so the hot, frequently-polled
        page-type tag does not false-share a cache line with the rest of the node's
        content; the root index itself moves from a separate \path{parking_lot} mutex
        onto the same \path{SmartCell}/\path{SmartGuard} machinery as ordinary blocks.
  \item \commit{e69ed87}/\commit{75453aa}/\commit{d15c693}/\commit{45cb048} (2026-05-26
        to 28, ``Addressing memory optimizations'') --- among other changes, guard
        lifetimes are changed from \path{'static} to an explicit \path{'a} (fixing a
        \path{mem::transmute}'d dangling reference for height-0 trees), and, tellingly,
        \commit{45cb048} \emph{reverts} an experiment that had range-query iterators hold
        a long-lived \path{BlockGuard} across the whole iteration: a \path{'static}
        lifetime guard held for an iterator's entire life is exactly the kind of
        long-lived, never-revalidated read that GC's reclaim decision does not know
        about. The revert restores a bare \path{BlockRef}, re-borrowed fresh on each
        step, keeping GC's liveness bookkeeping (\S\ref{sec:gc}) sound.
\end{itemize}

\section{Structural maintenance: splits, merges, and page capacity}
\label{sec:smo}

Structure-modification operations (SMOs) live in \path{src/mv_tree/smo.rs} and decide,
on the cold path, whether a page needs to split (overflow) or merge (underflow), and how.

\subsection{Capacity tuning}

Two independent tuning passes shrink the gap between ``nominally full'' and ``actually
lands cleanly on hardware pages'':

\begin{itemize}
  \item \commit{34555e2} (2026-06-20) --- \emph{``Updated overflow count, was set too
        eager. Its enough to trigger at $\le 1$ slot left for keys, not $\le 3$.''} The
        internal-page overflow trigger (\path{overflow_keys_count()}) moved from
        \path{max_keys() - 3} to \path{max_keys() - 1}, so a page is allowed to fill much
        closer to capacity before an SMO is even attempted.
  \item \commit{b29d347} (2026-08-02) --- \emph{``FAN\_OUT/NUM\_RECORDS: 125 $\to$ 123,
        so the real allocated block lands exactly on a 4KB page.''} The nominal record
        count was tuned down by two so that the \emph{actual} heap allocation
        (\path{OptCell<Block<..>>}, not \path{Block} alone --- the wrapping cell adds
        bytes) lands on a 4KB boundary: measured to move page-aligned allocations from
        25\% to 100\% of the time, with a 9--12\% reduction in dTLB load misses.
\end{itemize}

\subsection{Split: choosing where to cut}

\path{split()} decides between a \textbf{VERSION\_SPLIT} (compact obsolete/dead versions
out of one page in place) and a \textbf{KEY\_SPLIT} (cut the page in two at some key
boundary). Two real bugs surfaced here under concurrent TPC-C stress in July--August
2026:

\begin{itemize}
  \item \commit{2a2607e} (2026-08-02) --- \emph{``Fix split() off-by-one that let a
        full-but-protected leaf overflow.''} The VERSION\_SPLIT-vs-KEY\_SPLIT decision
        used a strict \path{survivor_count > capacity} test, so a page whose survivors
        landed at \emph{exactly} capacity wrongly took the ``just compact'' path,
        producing an already-100\%-full result with no room left for the pending write
        (a panic on a root-leaf tree, or an infinite retry loop otherwise). Fixed by
        \path{survivor_count >= capacity} (\path{smo.rs:1247}). This turned out to be the
        actual root cause of a stress-test failure originally suspected to be about
        hot-key version tearing (see \S\ref{sec:selfoverwrite}) --- the failing page in
        fact held several distinct keys, and simply had one transaction slower to commit
        than its neighbors.
  \item \commit{dc30111} (2026-08-02) --- \emph{``nearest\_key\_boundary: guarantee a
        capacity-fitting split when one exists.''} The boundary search used to combine
        two ``don't tear a key's version chain'' and ``fit within capacity'' constraints
        greedily and could miss a valid boundary that satisfied both (example from the
        commit: 31 records, capacity 16, the nearest non-tearing boundaries sit at 14 and
        17 --- both over capacity --- even though a capacity-fitting boundary exists
        further out). The fix searches in tiers, in order of preference: (1) no tear and
        fits, (2) fits but may tear, (3) no tear but may overflow, (4) fall back to the
        raw midpoint (\path{smo.rs:76-117}).
\end{itemize}

\subsection{Merge: choosing what to combine}

\begin{itemize}
  \item \commit{dcb5abe} (2026-08-02) --- \emph{``merge(): pick the emptier of two
        adjacent siblings, not always the right one.''} The candidate search now
        compares available room on both sides instead of defaulting to the right sibling.
  \item \commit{4c01360} (2026-08-02) --- \emph{``merge(): drop the pointless
        ts\_start ordering when combining two leaves.''} A \path{ts_start}-based re-sort
        risked reordering a live record ahead of a dead predecessor with a larger
        \path{ts_start}, corrupting the backward-scan (\path{rfind}) physical-order
        invariant relied on everywhere in \S\ref{sec:physical}.
  \item \commit{0659d34} (2026-08-02) --- a redundant \path{borrow_read().deref()}
        removed; a candidate-retry-on-failure idea was tried and reverted after it was
        found to spin without backoff and starve threads under contention.
  \item \commit{4bc78de} (2026-07-28) --- \emph{``Fix GC block reuse in merges; Now
        correctly calculates the live number.''} Merging two leaves based only on
        \path{active_count} could overflow \path{NUM_RECORDS}, because
        \path{record_survives_gc} (\S\ref{sec:physical}) also keeps GC-protected dead
        records that an active-count threshold alone does not see. Fixed by an explicit
        \path{leaf_merge_would_overflow} pre-check (\path{smo.rs:779-820}) that falls
        through to a \path{KeySplit} instead when the true survivor count would not fit.
\end{itemize}

\subsection{A week of hardening (2026-07-31 -- 08-02)}

The commits above were the tail end of a dense, one-week debugging arc chasing two real
TPC-C crashes: \commit{7432525} hardened \path{InternalPage} overflow checks without yet
fixing the crash; \commit{5c54658} fixed both crashes via ``GC-gated survivor
protection'' (a correct \path{record_survives_gc} distinguishing a genuinely-committed
delete from one whose abort-reversal still needs the predecessor physically intact) and
corrected \path{bulk_push} accounting (previously assuming every bulk-pushed record was
active); \commit{77a5689} reworked split/merge for a ``weakened fresh node rule''; and
\commit{fc8806b} reverted \path{on_overflow_node}/\path{on_underflow_node}/
\path{split_root}/\path{merge_root} back to \path{Result} + continue-on-success now that
the underlying bug was actually fixed, undoing a defensive workaround that had been
layered on top while the real cause was still unknown.

\section{The self-overwrite optimization}
\label{sec:selfoverwrite}

This is the trick at the heart of the user's original question: \emph{when a single,
still-uncommitted transaction updates the same key several times, overwrite the previous
update in place instead of appending a new multiversion entry.}

\subsection{The problem}

Recall from \S\ref{sec:physical} that every write mints one more physical
\path{RecordPoint} tuple, and that an \emph{uncommitted} tuple can be neither discarded
(the transaction might still commit) nor coalesced with anything (nothing has superseded
it yet, since nothing else can even see it). The old \path{DbTransaction::update} path
was literally ``delete old, push new'' on every call --- so if one open transaction wrote
the same key $N$ times before committing, it left $N$ physical, uncommitted, and
therefore un-reclaimable tuples for that one key in the page.

If $N$ grows large enough relative to a page's capacity, structural maintenance
(\S\ref{sec:smo}) is put in an impossible position: \path{split()}'s
\path{nearest_key_boundary} needs a key boundary to cut at, but a page dominated by one
key's own version chain has none --- the routine's own documentation
(\path{smo.rs:44-50}) states it falls back to tearing the chain across two disjoint-fence
siblings \emph{only} when the page is a single key's entire chain with no boundary at
all, silently orphaning the torn-off half from any fence-routed lookup. Symmetrically,
\path{merge()}'s combine-into-one-leaf path has to worry about overflowing
\path{NUM_RECORDS} on exactly the same accounting (\S\ref{sec:smo}). A concrete real-world
trigger: TPC-C's New-Order transaction pricing two order-lines for the same item lands
two updates on the very same \texttt{Stock} row inside one transaction.

\subsection{The fix}

Current code, \path{src/mv_db/transaction.rs:251-311} (\path{DbTransaction::update},
lightly abridged):

\begin{lstlisting}
match leaf_page.as_records_mut().iter_mut()
    .rfind(|r| r.key() == key && !r.version.insertion_stamp().is_invalid())
{
    Some(record) =>
        if tree.is_visible_stamp(self.worker_id, self.ts_start,
                                  record.version.insertion_stamp())
        {
            let stamp = TxStamp::new(self.worker_id, self.ts_start);
            tree.wal_log_write(stamp, |_| CRUDOperation::Update(key, payload.clone()));

            // Self-overwrite fast path: this record was already written
            // by *this* transaction -- nobody else can possibly be
            // relying on its current content, so mutate it in place.
            if record.version.insertion_stamp() == stamp {
                record.set_payload(payload);
                return CRUDOperationResult::Updated(stamp.ts_start());
            }

            // Otherwise: normal versioned update (delete old, push new).
            if !record.version.delete(stamp) { /* conflict */ }
            let len = leaf_page.len();
            leaf_page.push_uncommitted(
                RecordPoint::new(key, VersionInfo::new(stamp), payload), len);
            leaf_page.commit_delta(0, 1);
            self.written.borrow_mut().push((table, key));
            CRUDOperationResult::Updated(stamp.ts_start())
        } else { CRUDOperationResult::Conflict },
    None => CRUDOperationResult::ZeroAffected(KeyDoesNotExist),
}
\end{lstlisting}

\textbf{Detection} is a single equality check, not a heuristic: a \path{TxStamp} packs
\path{(worker\_id, ts\_start)} and is constant for the \emph{whole} transaction (one
stamp per transaction, not one per write --- \S\ref{sec:physical}). So ``was the newest
live version of this key written by me, in this very transaction'' reduces to
\path{record.version.insertion_stamp() == stamp}.

\textbf{Safety} follows directly from visibility: \path{is_visible_stamp} only returns
true for the same worker or a stamp already present in the commit log
(\S\ref{sec:osic}); neither applies to any other transaction while this one is still
open. There is therefore no possible reader anywhere relying on the current content of
this specific uncommitted tuple --- unlike a normal update, there is no
already-committed predecessor to protect, because nothing has committed yet. If the
found record instead belongs to another transaction, or is an older, already-committed
version, the code falls through to the ordinary delete-then-push path unchanged.

\subsection{Delete/reinsert cycles use the same bound}

Update self-overwrite alone did not cover a less common spelling of the same logical
operation. The first \path{delete(key)} followed by \path{insert(key, value)} in one
transaction must append a replacement tuple: the committed predecessor must remain
physically present so abort can undelete it. Before the follow-up fix, however, repeating
that pair in the same transaction appended another replacement on every cycle.

\path{insert_on_tree} now finds the newest non-invalid record before deciding to append.
If that record was both inserted and deleted by the current transaction --- its insertion
and deletion stamps equal the transaction's \path{TxStamp} --- insertion WAL-logs the new
logical value, clears the deletion mark, replaces the payload in place, and adjusts the
leaf's live/dead counters. It deliberately does not add another write-set entry: the
entry created for the first replacement already makes abort invalidate that tuple and
undelete the original predecessor. Thus $N$ delete/reinsert cycles create one replacement
tuple, commit exposes only the final payload, abort restores the pre-transaction value,
and WAL recovery replays the final logical state.

\subsection{The chronological story}

Two separate mechanisms are easy to conflate here; only the second is the trick this
section is about.

\paragraph{An older, unrelated mechanism: GC-driven update-in-place (2025--2026).}
\commit{bea1117} (2025-06-05) added an \emph{update-in-place} optimization to the
single-operation (non-transactional) CRUD dispatch path
(\path{src/mv_query/dispatch.rs}): if a record's insert version predates every live
snapshot --- i.e., GC has already determined nobody can still need the old
\emph{committed} value --- its payload is overwritten in place rather than versioned
again. \commit{16d1402}, \commit{7b1c466}, \commit{fad58a6}, \commit{7623f41}
(2026-02-02 to 02-04) incrementally wired and debugged the feature flag controlling this.
It is gated off whenever a WAL is attached (the fast path cannot be represented as a
single WAL record) and remains in the codebase today, but it is a GC/version-bloat
heuristic for single-op writes, unrelated to the transactional self-overwrite trick.

\paragraph{The actual fix (August 2026).}
\begin{enumerate}
  \item \commit{160186e} (2026-08-01) --- while adding concurrent TPC-C/YCSB/TPC-H stress
        tests, discovered that \path{DbTransaction::abort()} only reverted the
        \emph{newest} of several same-key writes, leaving older self-written versions
        permanently visible to that worker. Fixed \path{abort_write}/\path{apply_invalidate}
        to unwind the whole chain one physical version at a time. The commit message
        first names, but does not yet fix, the risk of a single hot key's version chain
        being torn across a same-key split.
  \item \commit{85ac1b0} (2026-08-01) --- \textbf{the fix itself}, as shown above. The
        commit message frames it exactly via the TPC-C Stock-row scenario. It also
        documents a \emph{reverted} companion idea --- lowering the leaf overflow
        trigger to buy compaction more headroom --- abandoned because it caused a hang
        under concurrent stress (a separate, pre-existing race left uninvestigated as
        out of scope).
  \item \commit{fbedd74} (2026-08-01) --- \emph{``Add minimal, ignored repro for the
        hot-key version-tearing limitation.''} A two-thread, \path{\#[ignore]}d repro
        isolating a \emph{different} contributor to the same class of pressure: a
        long-open snapshot keeping an old committed version of a hot key GC-protected
        while other transactions keep committing new versions of it --- explicitly no
        same-transaction repeats are involved. The test's own documentation states that
        the self-overwrite fast path ``cannot help here and doesn't change this repro's
        outcome --- it fixes a different, narrower contributor to the same class of page
        pressure, not this one.''
  \item \commit{2a2607e} (2026-08-02) --- tracing the repro further, the real cause
        turned out to be the split off-by-one from \S\ref{sec:smo}, not hot-key tearing
        at all: the failing leaf held several distinct keys. The repro file was renamed
        from \path{known_issue_hot_key_version_tearing_repro.rs} to
        \path{leaf_split_off_by_one_regression_tests.rs} to reflect the real cause.
\end{enumerate}

\subsection{Resulting bound and the separate TPC-C discrepancy}

$N$ updates, or $N$ delete/reinsert cycles, to one key by one transaction now cost at
most one unresolved replacement tuple. First-writer-wins prevents other transactions
from concurrently appending unresolved replacements for that same live key. After those
transactions resolve, ordinary GC compaction can discard obsolete committed versions;
pages containing multiple keys remain splittable at a key boundary. Consequently the
earlier proposed ``unsplittable hot-key chain protected by an old reader'' explanation
does not describe the current failure mode: the old reader stays on the historical
retired block rather than forcing all of its records into each replacement page.

The TPC-C stress invariant still permits a narrow 0.5\% difference between aggregate
\path{Stock} counters and inserted \path{OrderLine} rows. Temporarily making that
assertion exact reproduced a small accounting/update discrepancy in both GC modes. It
is tracked as a separate issue; neither the evidence nor the data structure semantics
support attributing it to same-key splitting or reader snapshot retention.

\section{Garbage collection: reclaiming dead blocks without contention}
\label{sec:gc}

\subsection{From a global mutex to sharded, lock-free reclaim}

\begin{itemize}
  \item \commit{acdc7a8} (2024-02-14) --- \emph{``Started GC.''} First implementation:
        a global \path{Mutex<RBTree<SnapShot>>} of active transactions and a global
        \path{Mutex<LinkedList<(*mut Version, BlockRef)>>} of dead pages; allocation
        pops the front of the dead-page list if it predates the oldest live snapshot.
        The commit message already names the problem that would dominate the next
        several months: \emph{``need to consider cases where merge might become
        impossible due to re-used pages\ldots locks must be left out, or possible
        deadlocks appear!''}
  \item \commit{2e3ad21} (2024-02-22) --- allocation switches to \path{try_lock()} on
        both structures (bail out rather than block), avoiding a blocking dependency
        between allocation and the block-latching path.
  \item \commit{95e149f} (2024-02-29) --- \emph{``Exchange GC ActiveTx with Concurrent
        B+Tree instead.''} The active-transaction index moves from a mutex-guarded
        red-black tree to the project's own OLC-locked concurrent B+-tree
        (\texttt{cc\_bplustree}), dispatched via \path{PeekMin}/\path{Insert}
        operations instead of a globally exclusive lock. (This design is itself later
        superseded --- see below.)
  \item \commit{2bae83c}, \commit{e4fd950} (2024-03-04, 03-07) --- two rounds of memory
        leak fixes in the reuse/dead-page path.
\end{itemize}

\textbf{Current end state} (\path{src/mv_gc/block_tracer.rs},
\path{src/mv_gc/tracker_handle.rs}): liveness is driven directly off
\path{TxContext::live_min_snapshot()} (\S\ref{sec:sharding}) rather than a separate
active-tx tree, and reclaim is driven off a \emph{worker-sharded} \path{SkipMap} rather
than a single list or B-tree --- both changes are part of the sharding story in
\S\ref{sec:sharding}.

\subsection{Two more races found by ThreadSanitizer, and how each was closed structurally}

\begin{itemize}
  \item \commit{d0e0977} (2026-07-21) --- \emph{``Stabilize GC and liveness in
        nodes, now re-added the obsolete/retired mark explicitly.''} A separate
        \path{retired: AtomicBool} field is folded into the \path{RETIRED_FLAG_VERSION}
        bits of the SmartCell's own \path{cell_version} word (\S\ref{sec:cc}). Before
        this, a two-step ``mark retired, then let \path{Drop} unlock later'' pattern
        raced a concurrent reuse's \path{clear_retired()} against the dying guard's
        deferred \path{Drop}, which stored a value computed from lock-\emph{acquisition}
        -time state and could silently revert a freshly reused cell's \path{cell_version}
        --- confirmed empirically as hangs/livelocks under GC plus heavy concurrency.
  \item \commit{ecc9900} (2026-07-27) --- \emph{``Fix issues GC related and stabilize
        SMO with racey stamps.''} Two independent bugs:
        \begin{enumerate}
          \item A \path{DeadPageKey} collision: the dead-block index was keyed by a bare
                \path{Version}, but a merge's two dying siblings deliberately share one
                death version (the replacement block's birth version); \path{SkipMap::insert}
                on a duplicate key silently displaces the existing entry, permanently
                leaking one of the two retired blocks. Fixed by keying on
                \path{(Version, block\_address)} instead.
          \item A genuine torn-read race in \path{VersionInfo}: \path{insert_stamp}/
                \path{delete_stamp} were plain, non-atomic \path{TxStamp} fields mutated
                by \path{delete}/\path{invalidate} while lock-free OLC readers read the
                same word with no lock and no re-validation --- confirmed via
                ThreadSanitizer. Fixed by backing both stamps with \path{AtomicU64} and
                \path{Relaxed} load/store (sufficient because the hazard was torn reads,
                not ordering).
        \end{enumerate}
        The same commit adds \path{nearest_key_boundary} (\S\ref{sec:smo}) to prevent a
        merge/split cut from tearing a single key's version chain across two
        disjoint-fence siblings.
  \item \commit{fe03d90} (2026-07-28) --- \emph{``Remove version\_region atomic need;
        Changed SMO Liveness, now only checks overlap without actual need to obsolete
        mark.''} The capstone of this whole thread: \path{InternalPage::version_region}
        changes from an array of \path{AtomicVersion} to plain \path{Version}, and the
        \path{mark_version_obsolete} mechanism --- an in-place \path{fetch_or} on an
        already-published, concurrently lock-free-read slot, itself confirmed via
        ThreadSanitizer as a genuine race against range-query iteration --- is deleted
        outright. It is replaced by \path{InternalPage::live_mask()}: slot $i$ is dead
        \emph{iff} some later slot $j > i$'s key interval overlaps slot $i$'s --- an
        $O(\mathrm{FAN\_OUT}^2)$ pure-integer check that runs only on the SMO cold path,
        exploiting the invariant that live sibling intervals are always pairwise
        disjoint and supersession is append-only. No writer ever mutates a slot after
        publication, so the race does not merely get fenced correctly --- it stops being
        possible.
\end{itemize}

\textbf{Recurring pattern.} Nearly every ``racey''/``stabilize'' fix in this section
follows the same shape: something was mutated \emph{in place} on an already-published,
lock-free-readable slot (an obsolete-mark bit, a retired bit, a delete/invalidate stamp),
requiring ever more careful atomic/fence discipline to patch around it --- until the
design was eventually changed so that nothing is ever mutated after publication at all.
This mirrors the fence-removal trajectory of \S\ref{sec:cc}: converging from ``atomics
and fences bolted on to patch races'' toward ``structure the data so the race is
structurally impossible.''

\section{Root indexing strategies}
\label{sec:root}

A reader must first find ``the root as of snapshot $S$'' before it can traverse anything.
\path{RootIndexType} (\path{src/mv_root/index_root.rs:28-55}) selects one of four
structures for this, all keyed by version (Table~\ref{tab:root}).

\begin{table}[h]
\centering
\begin{tabular}{@{}p{2.7cm}p{4.3cm}p{3.2cm}p{4.3cm}@{}}
\toprule
\textbf{Strategy} & \textbf{Structure} & \textbf{Lookup} & \textbf{Notes} \\
\midrule
Vanilla & \path{LinkedList<Root<..>>} & $O(n)$ linear \path{rfind} & Simplest baseline. \\
SK (SkipList) & \path{crossbeam_skiplist::SkipMap} & $O(\log n)$ expected,
  \path{range(..=si).last()} & Lock-free concurrent inserts/lookups. \\
DexaBTree & external OLC-locked \path{BPlusTree}, keyed by \path{SnapShot} &
  $O(\log n)$ via \path{PeekMax}/\path{Pred}/\path{Insert} CRUD dispatch & Reuses a full
  concurrent B+-tree just to index root pointers by version. \\
FrugalList (default) & lock-free singly-linked list with a probabilistic
  skip-list ``tower'' (\path{v\_ridgy} links, coin-toss $p=0.5$) & expected sub-linear via
  tower, falling back to the flat chain & Appends are a single \path{AtomicPtr} store
  (no CAS loop needed, since appends are already serialized by the enclosing SmartCell's
  write latch); much smaller per-node footprint than a full skip-list tower. \\
\bottomrule
\end{tabular}
\caption{Root-index strategies and their tradeoffs.}
\label{tab:root}
\end{table}

These are \emph{not} partitioned by version range --- each strategy is one single,
tree-wide structure holding every historical root, chosen once at construction. They
optimize a different axis than sharding (\S\ref{sec:sharding}): the cost of finding the
root pointer itself, a potential single hot spot every reader must consult, trading
lookup complexity against per-append cost and memory overhead.

\section{Sharding as a system-wide pattern}
\label{sec:sharding}

The word ``sharding'' shows up in this codebase in four genuinely distinct places. None
of them is classic horizontal key-range partitioning of one index; each instead
partitions a specific hot, shared piece of \emph{bookkeeping state} by worker (or, in one
case, splits the index space itself by table).

\subsection{Tables as shards: one tree per table}

\commit{5f92e9f} (2026-07-15) --- \emph{``Introduce Tables per indexes, not encoded into
a single tree.''} Before this commit, multiple logical tables were multiplexed into a
single tree. \path{Database} (\path{src/mv_db/database.rs:89-115}) now holds
\path{tables: ArcSwap<TableList<..>>}, a list of \path{TableEntry\{name, tree:
Arc<MVBTSt<..>>\}}: each table gets its own physically separate tree, its own root
index, and --- importantly for \S\ref{sec:gc} --- its own dead-block tracker
(\path{TrackerHandleSt} is genuinely per-table). Concurrent transactions touching
different tables therefore no longer contend on one tree's root/internal-page latches or
one dead-block skip list. What \emph{must} stay global for cross-table snapshot-isolated
transactions --- the clock, commit-log pruning, snapshot liveness --- moves into one
shared, non-generic \path{TxContext} introduced in the same commit.

\subsection{GC dead-block tracker: own-shard-first, round-robin steal}

\path{src/mv_gc/block_tracer.rs:49-62} --- \path{BlockTrace} holds one \path{SkipMap}
shard per CPU (\path{shard\_for(worker\_id) = worker\_id \% shards.len()}). Registration
of a died page (\path{register_died_page}) only ever inserts into the dying worker's own
shard. Reclaim (\path{try_reclaim}) pops from the caller's own shard first; only if that
shard is empty (or not old enough) does it fall back to scanning other shards, starting
from a shared round-robin cursor so no single remote shard is favored repeatedly.

\textbf{Why:} death versions all come from one global clock, so before sharding,
concurrent SMOs from every worker inserted into a \emph{single} skip list clustering near
its current maximum key --- an always-contended tail. Since a reclaimer only needs
\emph{a} reclaimable block, not the globally oldest one, trading global pop-order
precision for per-worker insert locality is a sound tradeoff. \commit{f9fc960}
(2026-07-23) is the exact patch that added the ``try own shard first'' fast path (before
it, \path{try_reclaim} always scanned round-robin from a shared cursor even when the
calling worker's own shard already had a usable entry).

A subtlety worth flagging precisely: \path{mv_wal/writer.rs}'s documentation uses
``shard'' language for its \path{hardened} watermark field, and \path{block_tracer.rs}
cites it as design precedent. In the current code, however, there is exactly \emph{one}
\path{WalWriter} per \path{Database} with a single \path{hardened: Arc<AtomicVersion>}
--- no \path{Vec<WalWriter>} or multi-shard aggregation actually exists today. The
``shard'' terminology there is vestigial documentation carried over from an earlier
design, not a live mechanism.

\subsection{\texttt{TxContext}: cache-padded, per-worker atomic slots}

\path{src/mv_sync/tx_context.rs:44-163} holds three parallel vectors, one slot per
\path{WorkerId}: \path{live_tx} (each worker's currently published outermost snapshot),
\path{live_tx_depth} (reentrancy depth, so only the outermost \path{begin_snapshot} call
touches \path{live_tx}), and \path{in_flight_bound} (a conservative lower bound published
while a worker is mid-registration, i.e.\ has drawn a \path{ts\_start} but not yet
inserted into \path{live\_tx}). Each worker only ever writes its own slot;
\path{free_block}/\path{live_min_snapshot} scan all slots with \path{Acquire} loads,
$O(\mathrm{max\_workers})$, with no cross-slot synchronization required.

This design directly replaces two earlier shared, globally-contended structures, and the
commit history records the profiling numbers that motivated the replacement:

\begin{itemize}
  \item \path{live_tx} replaced a shared \path{SkipMap<SnapShot, AtomicUsize>}
        (\path{mv\_gc::query\_tracer}) that, per an in-code comment, was measured eating
        ``roughly a quarter to half of all CPU cycles'' on a saturated TPC-C run.
  \item \path{in_flight_bound} replaced a single global
        \path{registrations\_in\_flight: AtomicUsize} that, instrumented under a
        16-thread update-heavy YCSB workload, rejected 73.8\% of \path{free_block}
        attempts (versus 5.2\% at 2 threads).
\end{itemize}

\commit{eac6e40} (2026-08-02) then wrapped all three \path{Vec<Atomic*>}s in
\path{crossbeam_utils::CachePadded}: previously, up to 8 adjacent workers' slots could
share one 64-byte cache line --- false sharing that directly contradicted those fields'
own design-intent comments (``no shared cache line with any other worker's slot''). This
is a pure layout change (\path{CachePadded<T>} derefs transparently) targeted at
64-core/128-thread scale, where cross-CCD coherence traffic is expensive.

\subsection{Thread-local worker identity and snapshot cache}

\path{src/mv_sync/worker.rs} hands out a \path{WorkerId} once per (tree, thread),
cached in a thread-local keyed by a process-wide \path{tree\_uid}. A second thread-local,
\path{SNAPSHOT\_CACHE}, gives each worker its own \path{SnapshotCache} --- the OSIC
paper's ``thread-local snapshot cache'' (\S\ref{sec:osic}) --- rather than a
\path{Vec} on the tree indexed by \path{WorkerId}: using actual thread-local storage
means no thread ever \emph{could} touch another worker's cache, a stronger guarantee than
``happens to always pass its own id.'' This is the assignment mechanism that feeds
\path{worker\_id} into both the \path{TxContext} slots above and the GC block-tracer
shards.

\section{OSIC integration: instant commit, lazily hardened}
\label{sec:osic}

\subsection{What OSIC is}

Ordered Snapshot Instant Commit \cite{alhomssi2023scalable} decouples \emph{visibility}
from \emph{durability}. Every transaction draws two timestamps from one global,
monotonically increasing atomic counter (the Global Logical Clock, GLC): \path{ts_start}
at begin, \path{ts_commit} at commit. Each of a fixed set of long-lived worker threads
keeps its own small, append-only commit log --- an ascending list of its own committed
timestamps. A write becomes visible to another worker's snapshot the moment its writer's
commit timestamp is appended to that worker's log: an $O(1)$, purely in-memory step ---
there is no waiting for the WAL to reach disk. Durability catches up asynchronously and
in batches (\emph{group commit}): writes and commit markers are appended to a WAL, and a
background thread periodically flushes/fsyncs everything accumulated so far in one I/O,
publishing a ``hardened'' watermark that callers can poll explicitly if they need a
durability guarantee.

\subsection{How cMVBT's commit path works today}

\begin{itemize}
  \item \path{src/mv_sync/clock.rs} --- \path{GlobalClock}: a single
        \path{AtomicVersion}, \path{fetch\_add(1, SeqCst)} mints both \path{ts\_start}
        and \path{ts\_commit}. There is no separate ``publish'' step: a drawn timestamp
        is immediately comparable.
  \item \path{src/mv_sync/commit_log.rs} --- \path{CommitLog}, one per worker, a
        \path{Mutex<Vec<Version>>}. \path{commit()} draws \path{ts\_commit} then pushes
        it under the lock; this ordering (the mutex serializes the append against any
        concurrent reader) is exactly what prevents a reader from observing a
        reserved-but-not-yet-inserted commit stamp. \path{commit\_pruned} additionally
        trims the log once it exceeds \path{max\_workers} entries still needed as
        \emph{some} live snapshot's lower commit bound (LCB).
  \item \path{src/mv_sync/visibility.rs} --- \path{is_visible}: a direct port of the
        paper's core visibility check. An invalid (aborted) stamp is never visible;
        same-worker writes are visible if \path{stamp.ts\_start() <= reader\_ts\_start};
        otherwise visibility is \path{LCB(stamp.worker\_id, reader\_ts\_start) >
        stamp.ts\_start()}, computed via each thread's own \path{SnapshotCache} memo.
  \item \path{src/mv_wal/writer.rs} --- \path{WalWriter}: a lock-free MPSC channel to one
        background flush thread. \path{log\_with\_stamp}/\path{log\_commit} just encode
        and enqueue, returning immediately (fire-and-forget, no per-op fsync wait).
        \path{flush\_loop} sleeps a short \path{GROUP\_COMMIT\_LINGER} (200\,$\mu$s)
        after the first message, then drains everything else non-blockingly, batches it
        into one \path{write\_all} + \path{sync\_data()}, and advances the
        \path{hardened} watermark to the maximum \path{ts\_start} of any \emph{Commit}
        message in that batch --- never based on a bare write.
  \item \path{src/mv_wal/record.rs} --- the wire format distinguishes
        \path{WalEntry::Write} from \path{WalEntry::Commit\{stamp, ts\_commit\}}: a
        commit marker is a separate append-only entry because a write is logged
        optimistically \emph{before} its transaction's fate is known.
  \item \path{src/mv_wal/recovery.rs} --- replay keeps only \path{Write}s with a matching
        \path{Commit} marker (commit-gated replay), and replays sorted by
        \path{(ts\_commit, file order)}, not by \path{ts\_start}, since commit order can
        differ from log order across concurrent workers.
  \item \path{src/mv_db/transaction.rs} --- \path{DbTransaction::commit}: appends to the
        committing worker's \path{CommitLog} (making every write from every table this
        transaction touched visible everywhere, instantly), then logs exactly one
        fire-and-forget WAL commit marker. \emph{This is the whole ``instant commit'':}
        one in-memory append, with nothing on the write/commit path waiting on disk I/O.
\end{itemize}

\textbf{Contrast with naive MVCC.} A naive ``lock, fsync, then publish'' design pays a
full flush round-trip, serially, on every commit. Here, visibility and durability are
fully decoupled: visibility costs one atomic append; durability is a separate, batched,
background concern, exposed only via an explicit
\path{wait\_wal\_hardened()}/\path{wal\_hardened\_version()} API for callers who
specifically need a durability checkpoint.

\subsection{Highlights and differences versus the paper}

\begin{table}[h]
\centering
\begin{tabular}{@{}p{4.3cm}p{4.3cm}p{5.9cm}@{}}
\toprule
\textbf{Paper concept (LeanStore/OSIC)} & \textbf{cMVBT realization} & \textbf{Notable difference} \\
\midrule
Global Logical Clock (GLC) & \path{GlobalClock}, single \path{fetch\_add} & Direct port; simplified from an earlier ad hoc thread-indexed watermark array (see below). \\
Per-worker commit log + LCB & \path{CommitLog} + \path{SnapshotCache} & Direct port, with an added \path{commit\_pruned} mode that trims the log even when block-reclaim GC is disabled, since an OSIC-style log otherwise grows unboundedly during long GC-off historic-query runs. \\
Early Lock Release (\S3.4 of the paper) & \emph{Removed} & The paper's pessimistic per-commit durability cap (blocking dependents on one's own flush) was replaced by a simpler per-writer three-state watermark advanced per \emph{batch}, with no caller blocking by default --- judged not worth the cost of a flush round-trip on every single operation. \\
Version storage & Versions live directly inside the multiversion B-tree's own pages (\S\ref{sec:physical}), not a separate version chain/undo log as in LeanStore's single-version core & OSIC's visibility/commit machinery is layered on top of an already-multiversion structure, so ``old version still visible'' is a property of the tree itself, not of a side chain that the commit protocol must also manage. \\
Group commit & \path{GROUP\_COMMIT\_LINGER} batching in \path{flush\_loop} & Direct port of the batching idea, tuned empirically (200\,$\mu$s) for this workload mix. \\
\bottomrule
\end{tabular}
\caption{OSIC as described in \cite{alhomssi2023scalable} versus its realization here.}
\end{table}

\subsection{Evolution: two eras}

\paragraph{Pre-OSIC era (Dec 2025 -- Jan 2026).} Before the paper's design was adopted
wholesale, the system used a single global commit-version counter with an explicit
publish step, plus a hand-rolled sharded thread-indexed ``committed'' watermark array ---
in effect, an imperfect reinvention of the CommitLog/GLC pair:
\begin{itemize}
  \item \commit{17b8f5f} (2025-12-17) --- \emph{``Changed commit from CAS to fetch\_add
        via SeqCst.''} Replaced a \path{compare\_exchange\_weak} retry loop with a bare
        \path{fetch\_add}, since commit numbering no longer needed a conditional check.
  \item \commit{0358c1c}/\commit{194f973}/\commit{87e2523} (2025-12-20/21) --- \emph{``Fix
        SI for update operations; divided into two atomic sequential operations, as
        described in the paper.''} The old update path interleaved commit-number
        allocation with the leaf mutation itself, using an admitted stopgap: a
        \path{delete(key, committed\_version + 1)} ``fixed visibility window'' fudge
        factor (the commit's own message calls this out: \emph{``For now, I use a fixed
        visibility window''}). One day later, this was replaced with the discipline the
        current code still follows: publish the new version first, \emph{then} mark the
        old one deleted --- closing the gap where an old version could go invisible
        before its replacement was guaranteed visible.
  \item \commit{047d094}, \commit{20396d7}, \commit{efdf7a6}, \commit{408becb},
        \commit{4981858}, \commit{67606e4}/\commit{11effc9}/\commit{7d43e8e},
        \commit{836a1db} (Jan 2026) --- a sequence of patches to a
        thread-id-indexed watermark array, repeatedly fixing edge cases around threads
        changing roles (an OLTP worker becoming a reader mid-run) or exiting, and a
        query tracer needing its own clock thread-id. The commit messages read as a
        slow discovery, by trial and error, of exactly the problem OSIC's per-worker
        commit-log design solves cleanly.
\end{itemize}

\paragraph{OSIC era (from \commit{714a82f}, 2026-07-10).} A large commit,
\emph{``OSIC/WAL and TPC-C/OLAP Benchmark,''} wholesale replaces the ad hoc clock/
watermark scheme with the paper's actual design: \path{GlobalClock}, per-worker
\path{CommitLog}, \path{WorkerRegistry} + \path{SnapshotCache}, and the entire
\path{mv\_wal} module, alongside new TPC-C/OLAP benchmark harnesses.

\paragraph{Post-OSIC refinement (Jul 13 -- 28, 2026).}
\begin{itemize}
  \item \commit{3e8b837} --- \emph{``Add group commit linger to collect and fsync logs
        more frequently.''} Added the 200\,$\mu$s linger delay described above, trading a
        little latency for far higher throughput per fsync.
  \item \commit{ab10660} --- \emph{``Make WAL worker non-blocking for callers and batch
        hardening.''} Removed the paper's Early Lock Release machinery (a
        \path{hardened} field plus \path{mark\_pending}/\path{mark\_resolved} letting a
        committing worker cap the global durability watermark at its own
        \path{ts\_commit - 1} until its WAL entry flushed, blocking dependents) in favor
        of a per-writer three-state watermark advanced per batch with no default
        blocking.
  \item \commit{73cc4a5}/\commit{a5e6b7a} --- \emph{``Add commit log truncate setting,
        turned on by default (gc = off + truncate\_commit\_log = on).''} Added
        \path{allow_historic_query(enabled)} (superseding an earlier,
        confusingly-named \path{truncate\_commit\_log(enabled)}), coupling and inverting
        the tradeoff explicitly: enabling historic queries disables both GC and
        commit-log truncation, since a run that needs to query old snapshots cannot
        prune the very log that lets it find them, while a normal OLTP run should prune
        aggressively to bound memory (documented cost without pruning: page faults from
        \path{Vec} growth showing up as $\sim$2.4\% of \path{commit}'s own profiled
        samples).
  \item \commit{ecc9900}, \commit{fe03d90} --- the same GC/SMO stabilization fixes
        described in \S\ref{sec:gc}, adjacent to (and interacting with) the OSIC-visible
        commit stream.
\end{itemize}

\paragraph{A lock-free WAL backend, and two concurrency bugs it took to land it (Aug 10,
2026).} Two independent lines of work attacked the same WAL append-path contention at
once, and had to be reconciled on merge.

\begin{itemize}
  \item \commit{4a1a50a} --- \emph{``Fix dead target-cpu=native, table-driven CRC32, shard
        WAL writer channel.''} Profiling (see \path{docs/oltp\_wal\_optimization.md}) found
        the single shared \path{crossbeam-channel} every worker's \path{WalWriter::enqueue}
        sent through was a genuine contention point (86\% of \path{Sender::send} samples
        sitting in its internal spin-wait). Fix: split into \path{NUM\_SHARDS = 4}
        independent (channel, flush-thread) pairs keyed by \path{worker\_id \% 4}, all
        still writing through one shared \path{Arc<Mutex<File>>} so the on-disk format and
        recovery are unaffected. Landed directly on \path{main} (not through this repo's
        usual review-by-conversation path), alongside a table-driven \path{crc32} (was
        22.6\% of self time on a write-heavy workload; a byte-at-a-time bit-loop) and a fix
        for \path{target-cpu=native} silently not applying (Cargo only reads
        \path{rustflags} from \path{.cargo/config.toml}, never the manifest's
        \path{[build]} section).
  \item \commit{ebdd789} --- \emph{``Add lock-free WAL backend, fix RESTART\_TRACE OOM
        leak, add fast WAL perf tests.''} A second, concurrent design for the same
        contention: \path{LockFreeWalWriter}/\path{LockFreeWalBackend}
        (\path{mv\_wal::lockfree\_writer}/\path{backend}) removes the channel and dedicated
        writer thread entirely --- every worker reserves its own byte range via
        \path{tail.fetch\_add} and calls \path{pwrite} directly, optionally batching its
        own records locally (\path{LocalBatch}) before flushing. \path{WalBackend} wraps
        both designs behind one enum so callers pick either at attach time. Benchmarking
        this (\path{tests/tpcc\_wal\_backend\_bench.rs}) surfaced a real bug, unrelated to
        either WAL design: \path{mv\_test::RESTART\_TRACE} (a write-restart diagnostic,
        documented ``off by default'') had been left \path{true} in checked-in code. Every
        OCC restart recorded a \path{String} key into a process-lifetime global map that's
        never cleared; the comparison loop calls \path{run\_tpcc()} 6 times in one process,
        so this accumulated across every iteration until the fastest config
        (\path{lockfree-batch64}, 16 terminals --- generating restarts fastest of all of
        them) grew past 12GB RSS and got OOM-killed, taking the whole terminal session down
        with it. Fixed by restoring the documented default and adding
        \path{reset\_restart\_trace()}, called at the start of every \path{run\_tpcc()}.
        Also added two always-on (not \path{\#[ignore]}d), deliberately tiny WAL perf smoke
        tests (\path{tpcc\_wal\_perf\_tests.rs}/\path{ycsb\_wal\_perf\_tests.rs}, under 20s
        combined, RSS growth asserted under a 20GB ceiling) so a regression like this one
        is caught on every \path{cargo test}, not only in the full-scale run that found it.
  \item \commit{9ced1a3} --- \emph{``Merge origin/main (WAL channel sharding + table-driven
        CRC32), fix cross-shard hardened\_version race.''} Merging \commit{4a1a50a} and
        \commit{ebdd789} exposed a second, more serious bug in the sharded
        \path{WalWriter} itself: \path{hardened\_version} published every shard's own
        confirmed watermark into \emph{one} value via \path{fetch\_max}, correct only if
        every shard's watermark were directly comparable --- it isn't, since
        \path{ts\_start} is one global sequence spread across shards unpredictably by
        \path{worker\_id \% 4}, not partitioned per shard. One busy shard's flush could
        advance the shared watermark past a \path{ts\_start} a slower shard hadn't flushed
        yet. \path{tree\_wal\_consistency\_tests.rs}'s
        \path{concurrent\_db\_transactions\_across\_tables\_match\_shared\_wal\_exactly}
        caught it directly: a quarter of the written records missing from the
        reconstructed WAL (\path{1200} vs. the expected \path{1800}) because
        \path{wait\_wal\_hardened} returned early. A first fix attempt (minimum of every
        shard's raw confirmed value) traded the false-positive for a hang: a shard whose
        own submissions never reach some \emph{other} shard's higher watermark would never
        satisfy a plain \path{min}, even once fully caught up on everything it actually
        had. The actual fix gives each shard a \path{submitted} watermark alongside
        \path{confirmed} (updated in \path{enqueue}, strictly before the message is sent,
        so it can only ever be a stale overestimate of in-flight work), and the aggregate
        now treats a shard as imposing no constraint once \path{confirmed >= submitted}
        (fully drained) rather than comparing raw values across shards. Verified: the
        failing test run 5$\times$ with no flakiness, the full suite 3$\times$ clean.
  \item \commit{efc3773} --- \emph{``Skip the discarded flush ticket on the WAL hot path,
        right-size framing buffers.''} A short \path{perf record} pass (15s TPC-C/YCSB-A,
        WAL on) run once the append path itself was no longer the bottleneck surfaced two
        purely allocator-side hotspots the earlier channel contention had been masking.
        (1)~\path{WalWriter::enqueue} built a fresh \path{bounded(1)} ack channel on
        \emph{every} call to return a flush ticket that every production call site
        (\path{mv\_sync::version\_handle}) always discarded --- split each logging method
        into a ticket-less default (no channel built at all) plus an explicit
        \path{*\_with\_ticket} variant for the handful of tests that genuinely wait on one
        record. (2)~every framing buffer's \path{Vec::with\_capacity} hint (24--44 bytes)
        was tuned for this module's own \path{u64}-payload tests, not real payloads ---
        \path{YcsbRow} (1000 bytes at this benchmark's field sizes) is roughly 25$\times$
        that, \path{TpccRow}'s \path{Customer}/\path{Stock} variants routinely 7--10$\times$
        over, forcing several grow-and-copy reallocations per WAL write. Added
        \path{WalPayload::wal\_encode\_size\_hint()} (exact for \path{u64}, and for
        \path{YcsbRow}/\path{TpccRow} via overrides mirroring their own \path{wal\_encode}
        field-for-field) and \path{record::entry\_size\_hint()}, and swapped every
        hardcoded constant in \path{writer.rs}/\path{lockfree\_writer.rs} for it. Re-measured
        on the same short runs: TPC-C tpmC +12.6\%, YCSB-A ops/sec +25.3\%, and the
        synthetic \path{u64}-payload \path{WalWriter} microbenchmark at 16 threads +79\%
        (largest gain here since, against an 8-byte payload, the discarded ticket's channel
        allocation was close to the entire per-write cost). Not fully closed: YCSB's
        page-fault/allocator chain shrank but didn't disappear --- a fresh $\sim$1000-byte
        \path{Vec} per write is still a fresh allocation per write; a true cross-thread
        buffer pool was flagged as a follow-up, not attempted.
\end{itemize}

\subsection{Known, documented limitations}

\begin{itemize}
  \item \textbf{Unbounded CommitLog growth} under long GC-off historic-query runs: page
        faults from \path{Vec} growth are directly attributed to \path{commit}'s own call
        stack. A chunked/segmented log was considered but explicitly deferred (``a real
        rewrite of \path{prune}/\path{lcb\_index}'s cross-block search'').
  \item \path{MAX_WORKERS_CAP = 256} --- a fixed-size inline array backing
        \path{SnapshotCache}, generous but still a hard ceiling.
  \item \textbf{Abort ordering}: \path{DbTransaction::abort}/\path{Drop} must revert
        writes in strict LIFO order, because reverting exposes each key individually to
        concurrent readers mid-abort --- a documented but delicate invariant, not a
        generically solved problem.
  \item \textbf{Self-overwrite leaves one physical entry, not zero}: writing the same key
        $N$ times in one transaction still leaves one extra physical entry versus never
        writing it at all. First-writer-wins prevents another transaction from building
        a concurrent unresolved chain for that key. Repeated delete/reinsert cycles now
        also reuse the transaction-owned replacement tuple after the first reinsert,
        keeping the unresolved per-key survivor set bounded.
  \item The scope of \commit{ecc9900}'s torn-read fix is explicit in its own
        documentation: it closes one specific torn-read hazard, not every possible
        concurrent-mutation hazard on the lock-free OLC read path in general.
\end{itemize}

\section{Timeline}
\label{sec:timeline}

Table~\ref{tab:timeline} indexes the commits discussed in this document chronologically,
for cross-reference against \texttt{git log}.

\begin{longtable}{@{}p{2.1cm}p{1.8cm}p{9.5cm}@{}}
\toprule
\textbf{Commit} & \textbf{Date} & \textbf{Theme} \\
\midrule
\endhead
\commit{acdc7a8} & 2024-02-14 & GC started (\S\ref{sec:gc}) \\
\commit{95e149f} & 2024-02-29 & GC active-tx index moved to a concurrent B+-tree (\S\ref{sec:gc}) \\
\commit{d565b2d}--\commit{5c26abb} & 2024-01 -- 2024-12 & ORWC implemented, iterated, deprecated, purged (\S\ref{sec:cc}) \\
\commit{48c4ded}, \commit{a7bb29c}, \commit{a518eed} & 2024-05 -- 2025-04 & Fence removal / ordering tightening (\S\ref{sec:cc}) \\
\commit{0cd4fee}, \commit{fa82c8c}, \commit{d97de21} & 2024-03 & Early DBTracker / snapshot bookkeeping (\S\ref{sec:sharding}) \\
\commit{bea1117}, \commit{16d1402}--\commit{7623f41} & 2025-06 -- 2026-02 & GC-driven update-in-place (unrelated mechanism, \S\ref{sec:selfoverwrite}) \\
\commit{17b8f5f}, \commit{0358c1c}, \commit{194f973}/\commit{87e2523} & 2025-12 & Pre-OSIC commit/visibility fixes (\S\ref{sec:osic}) \\
\commit{047d094}--\commit{836a1db} & 2026-01 & Pre-OSIC thread-watermark patching (\S\ref{sec:osic}) \\
\commit{34555e2} & 2026-06-20 & Overflow-trigger tuning (\S\ref{sec:smo}) \\
\commit{2a2be93}, \commit{40b00fd} & 2026-06 & Further ordering relaxation (\S\ref{sec:cc}) \\
\commit{714a82f} & 2026-07-10 & OSIC/WAL adopted wholesale (\S\ref{sec:osic}) \\
\commit{3e8b837}, \commit{ab10660}, \commit{73cc4a5}/\commit{a5e6b7a} & 2026-07-13/15 & Group commit, ELR removal, log truncation (\S\ref{sec:osic}) \\
\commit{5f92e9f} & 2026-07-15 & Tables-as-shards (\S\ref{sec:sharding}) \\
\commit{f9fc960} & 2026-07-23 & GC own-shard-first + steal (\S\ref{sec:sharding}) \\
\commit{d0e0977} & 2026-07-21 & Retired-bit race fix (\S\ref{sec:gc}) \\
\commit{45c7c21}/\commit{fd92f38} & 2026-07-21 & Conservative child latch removed (\S\ref{sec:cc}) \\
\commit{7f2bd09} & 2026-05-21 & Cache-line isolation for node tag (\S\ref{sec:cc}) \\
\commit{77a5689}, \commit{5c54658}, \commit{fc8806b}, \commit{7432525} & 2026-07-31 & TPC-C crash hardening week begins (\S\ref{sec:smo}) \\
\commit{ecc9900} & 2026-07-27 & Racey-stamp + dead-page-key collision fixes (\S\ref{sec:gc}) \\
\commit{fe03d90} & 2026-07-28 & Obsolete-mark removed; overlap-based liveness (\S\ref{sec:gc}) \\
\commit{4bc78de} & 2026-07-28 & Merge-overflow accounting fix (\S\ref{sec:smo}) \\
\commit{160186e} & 2026-08-01 & Repeated-key abort-revert fix (\S\ref{sec:selfoverwrite}) \\
\commit{85ac1b0} & 2026-08-01 & \textbf{Self-overwrite fast path} (\S\ref{sec:selfoverwrite}) \\
\commit{fbedd74} & 2026-08-01 & Hot-key repro added (\S\ref{sec:selfoverwrite}) \\
\commit{2a2607e} & 2026-08-02 & Split off-by-one found and fixed (\S\ref{sec:smo}, \S\ref{sec:selfoverwrite}) \\
\commit{dc30111}, \commit{4c01360}, \commit{dcb5abe}, \commit{0659d34} & 2026-08-02 & Merge/boundary-search fixes (\S\ref{sec:smo}) \\
\commit{eac6e40} & 2026-08-02 & Cache-padded per-worker atomics (\S\ref{sec:sharding}) \\
\commit{b29d347} & 2026-08-02 & FAN\_OUT tuned for page alignment (\S\ref{sec:smo}) \\
\commit{4a1a50a} & 2026-08-09 & WAL channel sharding, table-driven CRC32 (\S\ref{sec:osic}) \\
\commit{ebdd789} & 2026-08-10 & Lock-free WAL backend; RESTART\_TRACE OOM fixed (\S\ref{sec:osic}) \\
\commit{9ced1a3} & 2026-08-10 & Cross-shard hardened\_version race fixed (\S\ref{sec:osic}) \\
\commit{efc3773} & 2026-08-10 & Discarded flush ticket + undersized WAL buffers fixed (\S\ref{sec:osic}) \\
\bottomrule
\caption{Chronological index of commits discussed in this document.}
\label{tab:timeline}
\end{longtable}

\section{Conclusion}

Read end to end, the history of this codebase tells a consistent story at every layer:
an initial design gets something working with coarse-grained locks, mutexes, or bolted-on
atomics/fences; profiling or ThreadSanitizer then reveals exactly where that coarseness
costs correctness (a race) or throughput (contention); and the fix is, more often than
not, not ``add a lock'' but ``restructure the data so the hazard cannot occur'' ---
append-only pages instead of in-place edits, interval-overlap-derived liveness instead of
a mutable obsolete bit, per-worker sharded bookkeeping instead of one shared structure,
instant in-memory commit instead of commit-time fsync. The self-overwrite trick in
\S\ref{sec:selfoverwrite} is a small, local instance of exactly this pattern: rather than
adding machinery to garbage-collect or specially-handle a hot key's pile of uncommitted
versions, the fix removes the pile-up from ever occurring in the first place.

\begin{thebibliography}{9}

\bibitem{becker1996asymptotically}
Becker, B., Gschwind, S., Ohler, T., Seeger, B., Widmayer, P.
\emph{An asymptotically optimal multiversion B-tree}.
The VLDB Journal, 5(4):264--275, 1996.

\bibitem{tonta2026multiversion}
Tonta, A., Seeger, B., Soisalon-Soininen, E.
\emph{Multiversion Concurrency Control for Multiversion B-Trees}.
arXiv preprint arXiv:2606.09133, 2026.

\bibitem{alhomssi2023scalable}
Alhomssi, A., Leis, V.
\emph{Scalable and Robust Snapshot Isolation for High-Performance Storage Engines}.
Proceedings of the VLDB Endowment, 16(6):1426--1438, 2023.
doi:10.14778/3583140.3583157.

\end{thebibliography}

\end{document}
